% Part: lambda-calculus
% Chapter: church-rosser
% Section: beta-eta-reduction

\documentclass[../../../include/open-logic-section]{subfiles}

\begin{document}

\olfileid{lam}{cr}{be}

\olsection{$\beta\eta$-হ্রাস}

$\beta\eta$-হ্রাসের ($\bered$) ক্ষেত্রে চার্চ--রসার ধর্ম সত্য।

% BN-SRC-300 TARGET-CORRECTION-BEGIN
\par\smallskip\noindent\textnormal{\small\textbf{উৎস-সংশোধন (BN-SRC-300):} হিমায়িত শৈলী-ফাইলে $\beredone$-এর মুদ্রণরূপ থাকলেও পাঠটি নিচের lemma-র জন্য কোনো গাণিতিক এক-ধাপ সম্বন্ধ সংজ্ঞায়িত করেনি; তারপর $\beta$-সংকোচনের সংকেত~$\bredone$-কে “$\eta$-conversion” বলে। বাংলা পাঠে $M\beredone M'$-কে স্পষ্টভাবে $M\bredone M'$ অথবা $M\eredone M'$-এর সংক্ষিপ্ত রূপ ধরা হয়েছে; প্রমাণে মূল $\eta$-ক্ষেত্রে সঠিক $\eredone$ এবং সমান্তরাল $\eta$-বিধি ব্যবহৃত হয়েছে, আর সামঞ্জস্য-ক্ষেত্রগুলি আরোহে বহন করা হয়েছে।} % BN-SRC-300 TARGET-CORRECTION-END
\begin{lem}\ollabel{lem:one-par}
  $M \beredone M'$ হলে $M \beredpar M'$।
\end{lem}

\begin{proof}
  $M \beredone M'$-এর !!{derivation}-এর উপর আরোহ প্রয়োগ করি।
  সামঞ্জস্য-ক্ষেত্রগুলিতে আরোহ-অনুমানকে সমান্তরাল বিমূর্তন বা প্রয়োগের
  সংশ্লিষ্ট বিধি দিয়ে বহন করি। যদি $M \eredone M'$ মূল
  $\eta$-সংকোচন দিয়ে হয় (অর্থাৎ,
  \olref[syn][eta]{defn:beredone}), তবে \olref[pbe]{thm:refl}-এর সঙ্গে
  সমান্তরাল ইটা-বিধি ব্যবহার করি। অন্য মূল ক্ষেত্রটি
  \olref[b]{lem:one-par}-এর মতো।
\end{proof}


\begin{lem}\ollabel{lem:par-red}
  $M \beredpar M'$ হলে $M \bered M'$।
\end{lem}

\begin{proof} $M \beredpar M'$-এর !!{derivation}-এর উপর আরোহ প্রয়োগ
  করি। প্রথম চারটি ক্ষেত্র আগের সমান্তরাল-থেকে-সাধারণ হ্রাসের প্রমাণের
  মতো।
% BN-SRC-302 TARGET-CORRECTION-BEGIN
\par\smallskip\noindent\textnormal{\small\textbf{উৎস-সংশোধন (BN-SRC-302):} হিমায়িত প্রমাণটি আরোহ ঘোষণা করে কেবল পঞ্চম বিধির ক্ষেত্রটি দেয় এবং প্রথম চারটি ক্ষেত্র সম্বন্ধে কিছুই বলে না। বাংলা প্রমাণে স্পষ্ট করা হয়েছে যে সেগুলি আগের সমান্তরাল $\beta$-হ্রাস প্রমাণের একই চারটি নির্মাণ, এখন $\bered$-এর ভিতরে।} % BN-SRC-302 TARGET-CORRECTION-END

  শেষ বিধিটি \olref[pbe]{defn:beredpar5} হলে কোনো $x$, $N$, $N'$-এর
  জন্য $M$ হলো $\lambd[x][Nx]$ এবং $M'$ হলো $N'$, যেখানে
  $x \notin \FV{N}$ এবং $N \beredpar N'$। অতএব প্রথমে
  $\eta$-সংকোচনে $\lambd[x][Nx]$-কে $N$-এ হ্রাস করতে পারি, তারপর
  $N \bered N'$ দেখানো $\beredone$ ধাপগুলির ধারা অনুসরণ করি; শেষ কথাটি
  আরোহ-অনুমান থেকে সত্য।
% BN-SRC-301 TARGET-CORRECTION-BEGIN
\par\smallskip\noindent\textnormal{\small\textbf{উৎস-সংশোধন (BN-SRC-301):} হিমায়িত পঞ্চম-বিধির ক্ষেত্রে কাঁচা $FV(N)$ লেখা হয়েছে। বাংলা পাঠে অধ্যায়টির সংজ্ঞায়িত সংকেত $\FV{N}$ ব্যবহার করা হয়েছে।} % BN-SRC-301 TARGET-CORRECTION-END
\end{proof}


\begin{lem}\ollabel{lem:str}
  $\beredpar$-কে ধারণকারী ক্ষুদ্রতম পরিযায়ী সম্বন্ধ হলো~$\bered$।
\end{lem}

\begin{proof}
  \olref[b]{lem:str}-এর মতো।
\end{proof}

\begin{thm}\ollabel{thm:cr}
  $\bered$ চার্চ--রসার ধর্ম পূরণ করে।
\end{thm}

\begin{proof}
  \olref[dap]{thm:str}, \olref[pbe]{thm:cr} এবং \olref{lem:str} থেকে।
\end{proof}
\end{document}
